\documentclass[11pt]{amsart}
\usepackage{amsmath,amssymb,amsthm,booktabs}
\usepackage[margin=1.15in]{geometry}
\usepackage[hidelinks]{hyperref}
\newtheorem{theorem}{Theorem}[section]
\newtheorem{proposition}[theorem]{Proposition}
\newtheorem{lemma}[theorem]{Lemma}
\newtheorem{corollary}[theorem]{Corollary}
\theoremstyle{remark}
\newtheorem{remark}[theorem]{Remark}
\newcommand{\PP}{\mathbb{P}}
\newcommand{\EE}{\mathbb{E}}
\newcommand{\RR}{\mathbb{R}}
\newcommand{\Tr}{\operatorname{Tr}}
\newcommand{\doi}[1]{\href{https://doi.org/#1}{doi:\nolinkurl{#1}}}
\newcommand{\arxiv}[1]{\href{https://arxiv.org/abs/#1}{arXiv:\nolinkurl{#1}}}
\setlength{\emergencystretch}{2.5em}

\title[Frozen edges of tridiagonal GUE dynamics]{Frozen edges of stationary tridiagonal dynamics with GUE spectra}
\author{Yitzchak Shmalo}
\address{Yitzchak Shmalo, Tel Aviv University and Heights Labs}
\email{yspennstate@gmail.com}
\date{November 2026}
\subjclass[2020]{60B20, 60J60, 15B52}
\keywords{Dyson Brownian motion, tridiagonal matrix models, Gaussian unitary ensemble, Airy line ensemble, operator
limits, reversible diffusions}

\raggedbottom
\begin{document}
\begin{abstract}
Vir\'ag asked whether matrix-valued Brownian motion has an operator limit. A natural candidate is a stationary
diffusion on tridiagonal matrices with the Dumitriu--Edelman law of the Gaussian unitary ensemble (GUE) at every time.
We show that the simplest such diffusion, with independent Ornstein--Uhlenbeck diagonal entries and independent radial
Ornstein--Uhlenbeck off-diagonal entries, has frozen edge eigenvalues on the time scale on which Dyson Brownian motion
moves. For each fixed $k$ the $k$ largest edge-scaled eigenvalues converge, uniformly on compact time intervals, to a
vector that does not depend on time, so they do not converge to the Airy line ensemble. The same holds for every
stationary reversible tridiagonal diffusion with this invariant law whose entry noise is bounded uniformly in the
dimension. An exact tridiagonal realization of Dyson Brownian motion with these marginals needs entry noise that grows
with the dimension. Vir\'ag's question remains open.
\end{abstract}
\maketitle

\section{Introduction}\label{sec:intro}

Vir\'ag's survey of operator limits of random matrices closes with a list of questions; Question~8 of
\cite{Vir14}, on page~25 of the arXiv version, asks: \emph{are there operator limits of matrix-valued (say Hermitian)
Brownian motion?} The static theory supplies the template. Dumitriu and Edelman \cite[Theorems~2.1 and~2.12]{DE02}
showed that the Gaussian beta ensemble is the spectrum of a random Jacobi matrix with independent Gaussian diagonal and
chi-distributed off-diagonal entries, and Ram\'irez, Rider and Vir\'ag \cite{RRV11} proved that, seen at the soft edge
on the scale $n^{-1/6}$, these matrices converge to the stochastic Airy operator; see \cite[Theorem~5.1]{Vir14}. In the
bulk the Sine$_\beta$ operator of Valk\'o and Vir\'ag \cite{VV17} plays the same role. A dynamic version would couple
the matrices over physical time so that the operators converge as processes and the spectrum moves like Dyson Brownian
motion.

Two features make the question delicate. First, a static law on tridiagonal matrices says nothing about two-time laws.
Second, the edge of the stochastic Airy operator is governed by a corner of the Jacobi matrix of size of order
$n^{1/3}$, which grows with $n$, so a fixed number of tridiagonal entries cannot determine the edge dynamics.
Edelman, Jeong and Nissim \cite[Theorem~3.1]{EJN24} proved that, for the tridiagonal model obtained from Dyson Brownian
motion on real symmetric matrices, each fixed finite set of entries converges to independent Ornstein--Uhlenbeck
processes with position-dependent rates, and they leave the operator limit as a problem \cite[Problem~4]{EJN24}.
Holcomb and Paquette \cite{HP17} gave tridiagonal realizations of Dyson Brownian motion driven jointly by the
eigenvalues and the spectral weights, and Allez and Guionnet \cite{AG13} constructed invariant matrix diffusions with
Dyson eigenvalue dynamics. The spectral side is well understood: the largest eigenvalues of stationary Dyson Brownian motion
converge at the edge to the Airy line ensemble \cite{CH14,Lan20,GXZ24}, and the general-beta line ensemble has been
characterized through its pole dynamics \cite{HZ26}. None of these results produces a convergent family of random
differential operators in physical time.

The simplest dynamics with the right marginals assigns to every Dumitriu--Edelman entry its own stationary diffusion:
an Ornstein--Uhlenbeck process for each diagonal entry and a radial Ornstein--Uhlenbeck process, of the dimension
dictated by the chi law, for each off-diagonal entry. These are the beta-Hermite processes of Li \cite{Li10} at
$\beta=2$, after a deterministic normalization; Li already observed that their eigenvalues do not follow Dyson's
equations, and Yabuoku \cite{Yab23} computed the eigenvalue equations of related independent-entry processes.
Akemann, Majumdar and P\"a{\ss}ler \cite{AMP26} started the same kind of process from the zero matrix, computed the
joint law of its eigenvalues at every time and studied it under stochastic resetting. Our main
result shows that at the edge these dynamics are far from Dyson's: on the time scale $n^{-1/3}$ on which the edge of
Dyson Brownian motion moves by order one, the edge of the tridiagonal process does not move at all.

\subsection{The process}\label{sec:process}
Throughout, a GUE matrix $G$ has independent entries $G_{ii}\sim N(0,1)$ and $G_{ij}=(U_{ij}+iV_{ij})/\sqrt2$ for
$i<j$, with $U_{ij},V_{ij}$ independent standard normal; its density is proportional to $\exp(-\Tr G^2/2)$ and its
spectral edge is near $2\sqrt n$. For $\beta=2$ the Dumitriu--Edelman law is the law of the real symmetric tridiagonal
matrix with independent entries
\begin{equation}\label{eq:DE}
a_i\sim N(0,1)\quad(1\le i\le n),\qquad b_i\sim\chi_{2(n-i)}/\sqrt2\quad(1\le i<n)
\end{equation}
on the diagonal and off the diagonal, and its eigenvalues have the joint law of the GUE eigenvalues
\cite[Theorem~2.12]{DE02}. We call \eqref{eq:DE} the Dumitriu--Edelman law; a real tridiagonal matrix does not, of
course, have the law of $G$ itself.

Let $a_1(t),\dots,a_n(t)$ be independent stationary real Ornstein--Uhlenbeck processes, $da=-a\,dt+\sqrt2\,dW$, and let
$Z_i(t)\in\RR^{2(n-i)}$, $1\le i<n$, be independent stationary vector Ornstein--Uhlenbeck processes whose coordinates are
independent copies of the same scalar process, independent of the $a_i$. Put
\begin{equation}\label{eq:process}
b_i(t)=|Z_i(t)|/\sqrt2,\qquad T_n(t)=\operatorname{tridiag}\bigl(b_{i-1}(t),a_i(t),b_i(t)\bigr).
\end{equation}
We call $T_n$ the tridiagonal Ornstein--Uhlenbeck process. Let $\lambda_{1,n}(t)>\dots>\lambda_{n,n}(t)$ be its
eigenvalues and define the edge-scaled eigenvalues
\begin{equation}\label{eq:X}
X_{j,n}(s)=n^{1/6}\bigl(\lambda_{j,n}(sn^{-1/3})-2\sqrt n\bigr).
\end{equation}

\begin{theorem}\label{thm:main}
The process $T_n$ is stationary, reversible and Markov, has continuous paths and simple spectrum at all times almost
surely, and has the Dumitriu--Edelman law \eqref{eq:DE} at every time. Let $k\ge1$ and $S>0$, put
$r_n=\sqrt{12\log n}$, and let $n$ satisfy $\lfloor n^{1/4}\rfloor\ge k$. Then
\begin{equation}\label{eq:main}
\EE\max_{1\le j\le k}\sup_{0\le s\le S}\bigl|X_{j,n}(s)-X_{j,n}(0)\bigr|^2
\le 60\,kS\left(2\Bigl(\frac{6r_n+\pi^2k^2+2}{\sqrt n}\Bigr)^{1/2}+4n^{-5}\right).
\end{equation}
The right-hand side is $O_k\bigl(S\,n^{-1/4}(\log n)^{1/4}\bigr)$, with an implied constant depending only on $k$.
\end{theorem}

\begin{corollary}\label{cor:frozen}
For each $k$ and $S$, the processes $(X_{1,n},\dots,X_{k,n})$ converge in law in $C([0,S],\RR^k)$ to the process that is
constant in time and equal at every time to a vector with the joint law of the $k$ largest points of the Airy point process.
For every $s>0$ the limit satisfies $X_1(s)=X_1(0)$ almost surely, while the top curve $\mathcal A_1$ of the Airy line
ensemble satisfies $\PP\{\mathcal A_1(s)=\mathcal A_1(0)\}=0$.
\end{corollary}

Section~\ref{sec:normal} fixes the clock: with the normalization above, the edge of the rate-one matrix
Ornstein--Uhlenbeck process at physical time $sn^{-1/3}$ corresponds to the Airy line ensemble at time $s$. The
corollary therefore says that the tridiagonal Ornstein--Uhlenbeck process cannot have the Airy line ensemble as its edge
limit on any time scale comparable to Dyson's.

The proof uses only reversibility and the invariant law, so it extends to a class of diffusions
(Corollary~\ref{cor:bounded}): if a stationary reversible diffusion on Jacobi matrices has the Dumitriu--Edelman
invariant law and its entry noise has a covariance operator bounded by a constant $K$ independent of $n$, its edge
freezes in the same way. Without reversibility there is still a pointwise restriction
(Proposition~\ref{prop:growth}): an exact tridiagonal realization of Dyson Brownian motion with Dumitriu--Edelman
marginals needs entry noise whose covariance operator has norm at least of order $n^{1/4}(\log n)^{-1/4}$. These
statements rule out a natural class of candidates. They do not rule out operator limits with unbounded or correlated
entry noise, nonreversible dynamics, other coordinate systems, or limits that are not tridiagonal at all, and they do
not answer Question~8 in either direction.

The proof of Theorem~\ref{thm:main} has three ingredients. A forward--backward martingale decomposition, in the spirit
of Lyons and Zheng \cite{LZ88}, bounds the supremum of an increment of a stationary reversible diffusion by its
\emph{carr\'e du champ} alone, with no control of the drift (Lemma~\ref{lem:maximal}). For the tridiagonal
Ornstein--Uhlenbeck process the carr\'e du champ of an eigenvalue is at most three times the inverse participation
ratio $\sum_i v_i^4$ of its eigenvector (Lemma~\ref{lem:noise}), whereas for the full matrix process it equals one.
Finally, an elementary variational argument shows that the top eigenvectors of the Dumitriu--Edelman matrix are spread
over at least $n^{1/4}$ coordinates (Lemma~\ref{lem:deloc}). The simulations of Section~\ref{sec:sims} suggest that
the expected inverse participation ratio is of order $n^{-1/3}$; any bound tending to zero suffices for the proof.

\subsection{Organization}
\begin{center}
\small
\begin{tabular}{ll}
\toprule
Section & Content\\
\midrule
\ref{sec:normal} & normalization of Dyson Brownian motion and the Airy clock\\
\ref{sec:maximal} & maximal inequality (Lemma~\ref{lem:maximal}) and eigenvalue noise (Lemma~\ref{lem:noise})\\
\ref{sec:proof} & delocalization (Lemma~\ref{lem:deloc}); proofs of Theorem~\ref{thm:main} and
Corollary~\ref{cor:frozen}\\
\ref{sec:extensions} & bounded noise (Corollary~\ref{cor:bounded}), exact realizations
(Proposition~\ref{prop:growth}), Lanczos model\\
\ref{sec:sims} & simulations of both processes\\
\bottomrule
\end{tabular}
\end{center}
In the simulations the two-time correlation of the largest edge-scaled eigenvalue of the tridiagonal
Ornstein--Uhlenbeck process increases with $n$ toward one, while for Hermitian Brownian motion it stays near $0.4$.

\section{Normalization and the clock}\label{sec:normal}

Let $H(t)$ be the stationary rate-one Hermitian Ornstein--Uhlenbeck process with the GUE law of
Section~\ref{sec:process} at each time: each diagonal entry is a scalar process $dx=-x\,dt+\sqrt2\,dW$, and the real
and imaginary parts of each off-diagonal entry are independent processes $dx=-x\,dt+dW$. If $B(u)$ is Hermitian
Brownian motion with $B(1)$ a GUE matrix, then $e^{-t}B(e^{2t})$ is this process. The eigenvalues of $H$ solve
\begin{equation}\label{eq:DBM}
d\lambda_j=\sqrt2\,dB_j+\Bigl(2\sum_{\ell\ne j}\frac{1}{\lambda_j-\lambda_\ell}-\lambda_j\Bigr)dt
\end{equation}
with independent Brownian motions $B_j$, which is Dyson Brownian motion at $\beta=2$ in this normalization; see
\cite[Section~4.3]{AGZ10}. The martingale part of $X_j=n^{1/6}(\lambda_j-2\sqrt n)$ has quadratic variation
$2n^{1/3}\,dt$, so on the clock $t=sn^{-1/3}$ it is $2\,ds$. In the normalization of Corwin and Hammond, the top curves
$\mathcal A_i$ of the Airy line ensemble are such that $2^{-1/2}(\mathcal A_i(s)-s^2)$ are locally standard Brownian
bridges \cite[equation~(4) and Theorem~3.1]{CH14}, so their quadratic variation is also $2\,ds$. Thus physical time
$sn^{-1/3}$ for $H$ corresponds to Airy time $s$, with no further constant. The same time scale appears from the
spectral side: the oscillator generator $-\partial_x^2+x^2/4-1/2$ of the scalar Ornstein--Uhlenbeck process, centred at
$x=2\sqrt n$ and rescaled by $x=2\sqrt n+n^{-1/6}z$, satisfies
$n^{-1/3}(-\partial_x^2+x^2/4-1/2-n)=-\partial_z^2+z+O(n^{-1/3})+O(n^{-2/3}z^2)$, the Airy operator.

\section{The noise of an eigenvalue}\label{sec:maximal}

This section bounds the movement of a stationary reversible diffusion by its carr\'e du champ and computes the carr\'e
du champ of an eigenvalue of the tridiagonal Ornstein--Uhlenbeck process.

\subsection{A maximal inequality for reversible diffusions}

Let $Y=(Y_t)_{t\ge0}$ be a stationary Ornstein--Uhlenbeck process on $\RR^N$ with independent rate-one coordinates,
$dY=-Y\,dt+\sqrt2\,dW$, with generator $L=\Delta-y\cdot\nabla$, invariant law the standard Gaussian measure $\gamma$,
and carr\'e du champ $\Gamma(f)=|\nabla f|^2$.

\begin{lemma}\label{lem:maximal}
For every Lipschitz function $f:\RR^N\to\RR$ and every $U>0$,
\[
\EE\sup_{0\le t\le U}\bigl|f(Y_t)-f(Y_0)\bigr|^2\le 20\,U\int\Gamma(f)\,d\gamma .
\]
\end{lemma}

\begin{proof}
First let $f$ be smooth with bounded derivatives of all orders. By It\^o's formula,
$M_t=f(Y_t)-f(Y_0)-\int_0^tLf(Y_r)\,dr$ is a martingale with $d\langle M\rangle_t=2\Gamma(f)(Y_t)\,dt$, so
$\EE M_U^2=2U\int\Gamma(f)\,d\gamma$ by stationarity. The Gaussian transition kernel of $Y$ is symmetric with respect to
$\gamma$, so the reversed process $\widehat Y_u=Y_{U-u}$, $0\le u\le U$, has the same law as $(Y_u)_{0\le u\le U}$. Let
$\widehat M$ be the corresponding martingale for $\widehat Y$, with respect to the filtration of $\widehat Y$. The drift
integrals cancel in the difference, and
\[
f(Y_t)-f(Y_0)=\tfrac12M_t-\tfrac12\bigl(\widehat M_U-\widehat M_{U-t}\bigr),\qquad 0\le t\le U.
\]
Hence $\sup_t|f(Y_t)-f(Y_0)|^2\le\frac12\sup_t|M_t|^2+2\sup_u|\widehat M_u|^2$. Doob's $L^2$ inequality bounds the
expectation of each supremum by $4\cdot 2U\int\Gamma(f)\,d\gamma$, which gives the constant $4+16=20$.

For Lipschitz $f$, let $f_\varepsilon=f*\rho_\varepsilon$ for a smooth compactly supported mollifier $\rho_\varepsilon$.
Then $f_\varepsilon$ is smooth with bounded derivatives, $f_\varepsilon\to f$ uniformly, and
$\nabla f_\varepsilon=(\nabla f)*\rho_\varepsilon\to\nabla f$ almost everywhere and boundedly, so
$\int\Gamma(f_\varepsilon)\,d\gamma\to\int\Gamma(f)\,d\gamma$. The suprema converge uniformly, and the inequality
passes to the limit.
\end{proof}

The proof uses neither a bound on the drift nor a spectral gap; only stationarity, reversibility and the martingale
calculus enter. The same argument applies to any stationary reversible diffusion for which these are available, with
$2\Gamma$ the density of the quadratic variation of $f(Y)$.

\subsection{The carr\'e du champ of an eigenvalue}\label{sec:noise}

We view $T_n$ as a function of the Gaussian coordinates $(a,Z)=(a_1,\dots,a_n,Z_1,\dots,Z_{n-1})$, which form a
stationary Ornstein--Uhlenbeck process in dimension $N=n+n(n-1)$ as in Section~\ref{sec:maximal}.

\begin{proof}[Proof of the structural statements of Theorem~\ref{thm:main}]
The process $(a,Z)$ is a stationary reversible Gaussian Markov process. The transition kernel of each $Z_i$ commutes
with rotations, so $|Z_i|$ is a Markov process and $(a,|Z_1|,\dots,|Z_{n-1}|)$ is a stationary reversible Markov
process. At each time its law is that of independent $a_i\sim N(0,1)$ and $|Z_i|\sim\chi_{2(n-i)}$, which is
\eqref{eq:DE} after dividing by $\sqrt2$. The smallest radial dimension is $2$, and the origin is polar for a
two-dimensional Ornstein--Uhlenbeck process, so all $b_i(t)$ are positive for all $t$ almost surely. A Jacobi matrix with
positive off-diagonal entries has simple spectrum, and the ordered eigenvalues depend continuously on the entries, so
the paths of $\lambda_{j,n}$ are continuous. By It\^o's formula the off-diagonal entries solve
\[
db_i=\Bigl(\frac{2(n-i)-1}{2b_i}-b_i\Bigr)dt+d\widetilde W_i
\]
with independent Brownian motions $\widetilde W_i$.
\end{proof}

\begin{lemma}\label{lem:noise}
Let $f=\lambda_{j,n}$, viewed as a function of $(a,Z)$. Then $f$ is $1$-Lipschitz, and at every point where the spectrum
is simple and $b_i\ne0$ for all $i$,
\[
\Gamma(\lambda_{j,n})=\sum_{i=1}^nv_i^4+2\sum_{i=1}^{n-1}v_i^2v_{i+1}^2\le3\sum_{i=1}^nv_i^4,
\]
where $v=(v_1,\dots,v_n)$ is a unit eigenvector of $T_n$ for $\lambda_{j,n}$.
\end{lemma}

\begin{proof}
The map $(a,Z)\mapsto T_n$ satisfies $\|T_n-T_n'\|_F^2=\sum_i(a_i-a_i')^2+2\sum_i(b_i-b_i')^2\le|(a,Z)-(a',Z')|^2$,
since $|b_i-b_i'|\le|Z_i-Z_i'|/\sqrt2$, and ordered eigenvalues are $1$-Lipschitz in the Frobenius norm. At a simple
eigenvalue, first-order perturbation gives $\partial\lambda_j/\partial a_i=v_i^2$ and
$\partial\lambda_j/\partial b_i=2v_iv_{i+1}$. Since $|\nabla_{Z_i}b_i|^2=1/2$,
\[
\Gamma(\lambda_j)=\sum_iv_i^4+\tfrac12\sum_i(2v_iv_{i+1})^2,
\]
and $2v_i^2v_{i+1}^2\le v_i^4+v_{i+1}^4$ gives the bound.
\end{proof}

For the full Hermitian process $H$ the same computation over all matrix coordinates gives $\Gamma(\lambda_j)=1$, in
agreement with the noise coefficient $\sqrt2$ in \eqref{eq:DBM}. Lemma~\ref{lem:noise} replaces this by the inverse
participation ratio of the eigenvector, which is small when the eigenvector is spread out.

\section{Proof of Theorem~\ref{thm:main} and Corollary~\ref{cor:frozen}}\label{sec:proof}
\subsection{Delocalization of the top eigenvectors}\label{sec:deloc}

Let $T$ have the Dumitriu--Edelman law \eqref{eq:DE}, with eigenvalues $\lambda_1>\dots>\lambda_n$ and unit
eigenvectors $v^{(j)}$, and let $r_n=\sqrt{12\log n}$. Consider the event
\begin{equation}\label{eq:En}
E_n=\Bigl\{\max_i|a_i|\le r_n,\ \max_{i<n}\bigl|b_i-\sqrt{n-i}\bigr|\le r_n\Bigr\}.
\end{equation}

\begin{lemma}\label{lem:deloc}
For $n\ge2$, $\PP(E_n^c)\le4n^{-5}$. If $\lfloor n^{1/4}\rfloor\ge j$, then on $E_n$
\[
\sum_i\bigl(v_i^{(j)}\bigr)^4\le2\Bigl(\frac{6r_n+\pi^2j^2+2}{\sqrt n+r_n}\Bigr)^{1/2},
\]
and consequently $\EE\sum_i(v_i^{(j)})^4\le2\bigl((6r_n+\pi^2j^2+2)/\sqrt n\bigr)^{1/2}+4n^{-5}$.
\end{lemma}

\begin{proof}
For $R\sim\chi_d$, $\EE e^{\theta R^2}=(1-2\theta)^{-d/2}$, and the Chernoff bound gives
$\PP(R^2\ge dy)\le\exp[-d(y-1-\log y)/2]$ for $y>1$ and $\PP(R^2\le dy)\le\exp[-d(y-1-\log y)/2]$ for $y<1$; see also
\cite[Lemma~1]{LM00}. Since $y-1-\log y\ge(\sqrt y-1)^2$, this gives $\PP(|R-\sqrt d|>u)\le2e^{-u^2/2}$. With
$d=2(n-i)$ and $u=\sqrt2\,r_n$, each $b_i$ leaves its window with probability at most $2e^{-r_n^2}=2n^{-12}$, and each
$a_i$ with probability at most $2e^{-r_n^2/2}=2n^{-6}$. A union bound gives $\PP(E_n^c)\le4n^{-5}$.

Assume $E_n$ and put $m=\lfloor n^{1/4}\rfloor\ge j$. Let $A$ be the leading $m\times m$ block of $T$ and let $B$ be the
$m\times m$ tridiagonal Toeplitz matrix with zero diagonal and off-diagonal entries $\sqrt n$. For $i<m$,
$0\le\sqrt n-\sqrt{n-i}\le m/\sqrt n$, so every entry of $A-B$ on the diagonal is at most $r_n$ and every off-diagonal
entry is at most $r_n+m/\sqrt n$ in absolute value. The row-sum bound gives $\|A-B\|\le3r_n+2m/\sqrt n\le3r_n+2$. The
eigenvalues of $B$ are $2\sqrt n\cos(\ell\pi/(m+1))$, and $2(1-\cos x)\le x^2$ with $(m+1)^2>\sqrt n$ gives
$2\sqrt n\cos(j\pi/(m+1))\ge2\sqrt n-\pi^2j^2$. By Cauchy interlacing and Weyl's inequality,
\begin{equation}\label{eq:lower}
\lambda_j\ge2\sqrt n-3r_n-\pi^2j^2-2 .
\end{equation}

Let $w_i=|v^{(j)}_i|$, $w_0=w_{n+1}=0$, $M=\sqrt n+r_n$ and $D=\sum_{i=0}^n(w_{i+1}-w_i)^2$. On $E_n$ we have
$|a_i|\le r_n$ and $0<b_i\le M$, so
\[
\lambda_j=\sum_ia_i(v^{(j)}_i)^2+2\sum_ib_iv^{(j)}_iv^{(j)}_{i+1}\le r_n+2M\sum_iw_iw_{i+1}=r_n+M(2-D),
\]
using $\sum_iw_i^2=1$. Together with \eqref{eq:lower}, $D\le(6r_n+\pi^2j^2+2)/M$. For every $i$,
\[
w_i^2=\sum_{\ell<i}(w_{\ell+1}-w_\ell)(w_{\ell+1}+w_\ell)\le D^{1/2}\Bigl(\sum_\ell(w_{\ell+1}+w_\ell)^2\Bigr)^{1/2}\le2D^{1/2},
\]
and $\sum_iw_i^4\le\max_iw_i^2$. Off $E_n$ the inverse participation ratio is at most one.
\end{proof}

An eigenvector spread evenly over a corner of size $n^{1/3}$ would have inverse participation ratio of order $n^{-1/3}$,
and the simulated values of $\EE\,\Gamma(\lambda_1)$ in Section~\ref{sec:sims} fall at about this rate. The bound of
Lemma~\ref{lem:deloc} is weaker, and it is enough for Theorem~\ref{thm:main}.

\subsection{Conclusion of the proofs}

Apply Lemma~\ref{lem:maximal} to $f=n^{1/6}\lambda_{j,n}$, which is Lipschitz by Lemma~\ref{lem:noise}, with
$U=Sn^{-1/3}$. The spectrum is simple and all $b_i$ are positive $\gamma$-almost everywhere, so Lemma~\ref{lem:noise}
identifies $\Gamma(f)$ almost everywhere and
\[
\EE\sup_{0\le s\le S}|X_{j,n}(s)-X_{j,n}(0)|^2\le20\,Sn^{-1/3}\cdot3n^{1/3}\,\EE\sum_i\bigl(v_i^{(j)}\bigr)^4
=60\,S\,\EE\sum_i\bigl(v_i^{(j)}\bigr)^4,
\]
the expectation taken under the Dumitriu--Edelman law. Bounding the maximum over $j\le k$ by the sum and applying
Lemma~\ref{lem:deloc} with $j\le k$ proves \eqref{eq:main}.

For Corollary~\ref{cor:frozen}, the vector $(X_{1,n}(0),\dots,X_{k,n}(0))$ has the law of the $k$ largest edge-scaled GUE
eigenvalues, which converges to the law of the $k$ largest points of the Airy point process \cite{TW94}. The map sending a
vector to the constant path is continuous from $\RR^k$ to $C([0,S],\RR^k)$, and by \eqref{eq:main} the uniform distance
between $(X_{j,n})_{j\le k}$ and the constant path at $(X_{j,n}(0))_{j\le k}$ tends to zero in probability. This proves
the convergence.

For the last statement, fix $s>0$ and an open interval $I\ni s$ whose closure does not contain $0$. By the Brownian Gibbs
property \cite[Theorem~3.1]{CH14}, conditionally on the curves outside $I\times\{1\}$, the shifted top curve
$2^{-1/2}(\mathcal A_1(x)-x^2)$ on $I$ is a Brownian bridge between its endpoint values conditioned to stay above the
second curve. The conditioning event has positive probability, since the endpoints lie strictly above the continuous
lower curve, so the conditional law is absolutely continuous with respect to the Brownian bridge law. The value of the
bridge at the interior point $s$ has a continuous distribution, whereas $\mathcal A_1(0)$ is fixed by the conditioning.
Hence $\PP\{\mathcal A_1(s)=\mathcal A_1(0)\}=0$. \qed

By stationarity, $\operatorname{Cov}(\mathcal A_1(0),\mathcal A_1(s))=\operatorname{Var}(\mathcal A_1(0))-\frac12
\EE(\mathcal A_1(s)-\mathcal A_1(0))^2$ is strictly smaller than the variance of the Tracy--Widom law, whereas the frozen
limit has correlation one at every lag. Section~\ref{sec:sims} compares the two numerically.

\section{Bounded noise and exact realizations}\label{sec:extensions}

The proof of Theorem~\ref{thm:main} used two facts about the dynamics: reversibility, through Lemma~\ref{lem:maximal},
and a bound on the carr\'e du champ of an eigenvalue in terms of the derivatives in the entries. Write $c=(a,b)$ for the
entries of a Jacobi matrix with positive off-diagonal entries.

\begin{corollary}\label{cor:bounded}
Let $c(t)$ be a stationary reversible diffusion on Jacobi matrices with the Dumitriu--Edelman invariant law
\eqref{eq:DE}, for which the forward and backward martingale calculus of Lemma~\ref{lem:maximal} holds for Lipschitz
functions of the entries. Suppose its carr\'e du champ satisfies
\begin{equation}\label{eq:bounded}
\Gamma_n(f)\le K\Bigl(\sum_i|\partial_{a_i}f|^2+\sum_i|\partial_{b_i}f|^2\Bigr)
\end{equation}
with $K$ independent of $n$. Then the conclusion of Theorem~\ref{thm:main} holds with $60$ replaced by $100K$.
\end{corollary}

\begin{proof}
The forward--backward argument gives Lemma~\ref{lem:maximal} with $\Gamma_n$ in place of $\Gamma$. By the derivative
formulas in Lemma~\ref{lem:noise},
$\Gamma_n(\lambda_j)\le K\bigl(\sum_iv_i^4+4\sum_iv_i^2v_{i+1}^2\bigr)\le5K\sum_iv_i^4$. Lemma~\ref{lem:deloc}
concerns only the invariant law and applies unchanged.
\end{proof}

Condition \eqref{eq:bounded} is a bound on the covariance operator of the entry noise in the coordinates $(a,b)$; bounded
individual covariances do not imply it. The assumption on the invariant law matters as well. The diagonal matrix of the
ordered eigenvalues of stationary Dyson Brownian motion is a reversible continuous diffusion on (degenerate) tridiagonal
matrices with bounded noise and GUE eigenvalues at every time, and its edge moves; its entry law is not \eqref{eq:DE}, and
Lemma~\ref{lem:deloc} fails for it. If $K=K_n$ depends on $n$, the conclusion persists as long as
$K_nn^{-1/4}(\log n)^{1/4}\to0$.

Reversibility is not needed for a pointwise restriction on exact realizations of \eqref{eq:DBM}.

\begin{proposition}\label{prop:growth}
Let $c(t)$ be a stationary continuous semimartingale on Jacobi matrices with positive off-diagonal entries and the
Dumitriu--Edelman law at each time, with quadratic covariation $d\langle c\rangle_t=2C_n(c_t)\,dt$, and suppose its
eigenvalues solve \eqref{eq:DBM}. Then for almost every $t$, almost surely,
\[
\|C_n(c_t)\|_{\mathrm{op}}\ge\Bigl(5\sum_i\bigl(v_i^{(j)}(t)\bigr)^4\Bigr)^{-1}\qquad(1\le j\le n).
\]
In particular $\|C_n(c_t)\|_{\mathrm{op}}\to\infty$ in probability, and on the event $E_n$ of \eqref{eq:En} it is at least
$\frac1{10}\bigl(\sqrt n/(6r_n+\pi^2+2)\bigr)^{1/2}$, which is of order $n^{1/4}(\log n)^{-1/4}$.
\end{proposition}

\begin{proof}
By It\^o's formula the martingale part of $\lambda_j(c_t)$ has quadratic variation
$2\nabla\lambda_j^{\mathsf T}C_n\nabla\lambda_j\,dt$, and by \eqref{eq:DBM} it is $2\,dt$. Hence
$1=\nabla\lambda_j^{\mathsf T}C_n\nabla\lambda_j\le\|C_n\|_{\mathrm{op}}|\nabla\lambda_j|^2$, and
$|\nabla\lambda_j|^2=\sum_iv_i^4+4\sum_iv_i^2v_{i+1}^2\le5\sum_iv_i^4$. The remaining statements follow from
Lemma~\ref{lem:deloc} with $j=1$.
\end{proof}

The rates in the fixed-corner limits of \cite{EJN24} therefore cannot be extended to a tridiagonal model of Dyson
Brownian motion with entry noise bounded uniformly in $n$.

\begin{remark}\label{rem:lanczos}
The tridiagonal model of Dyson Brownian motion obtained by Lanczos tridiagonalization is Markov, although it has no
simple coefficient equation. Let $\mathcal C$ be the open set of Hermitian matrices for which $e_1$ is a cyclic vector,
and let $J(H)$ be the Jacobi matrix with positive off-diagonal entries produced by the Lanczos recursion from $e_1$,
with $H=Q(H)J(H)Q(H)^*$ and $Q(H)e_1=e_1$. The recursion uses products, inner products and divisions by positive
residual norms, so $J$ and $Q$ are continuous on $\mathcal C$. With $K=\{U\in U(n):Ue_1=e_1\}$, one has
$J(UHU^*)=J(H)$ for $U\in K$, and conversely $J(H)=J(H')$ implies $H'=UHU^*$ with $U=Q(H')Q(H)^*\in K$, so the fibres of
$J$ are exactly the $K$-orbits. Let $P_t$ be the transition kernel of Hermitian Brownian motion or of the
Ornstein--Uhlenbeck process $H$; it commutes with conjugation by unitary matrices, so $\int g(J(H'))\,P_t(H,dH')$ is
constant on fibres. Define
\[
\overline P_tg(j)=\int g\bigl(J(H')\bigr)\,P_t(j,dH'),
\]
viewing $j$ as a Hermitian matrix. For $t>0$ the measure $P_t(H,\cdot)$ has a density, and the complement of
$\mathcal C$ is an algebraic set of Lebesgue measure zero, so any Borel extension of $J$ off $\mathcal C$ leaves
$\overline P_t$ unchanged. Then $\EE[g(J(H_{s+t}))\mid\mathcal F^H_s]=\overline P_tg(J(H_s))$ almost surely for each
deterministic $s$ and $t>0$, and conditioning on the smaller filtration of $J(H)$ gives the Markov property at
deterministic times; the Chapman--Kolmogorov identity follows from that of $P_t$. This is a statement at deterministic
times; it does not by itself give a version of $J(H_t)$ defined and continuous at all times simultaneously, which would
require showing that the path avoids the complement of $\mathcal C$. Proposition~\ref{prop:growth} applies to this model
whenever it is a semimartingale in the entries.
\end{remark}

\section{Simulations}\label{sec:sims}

We simulated both processes on the clock $t=sn^{-1/3}$ at $s=0,\frac14,1,3$, recording the two largest edge-scaled
eigenvalues. The dense model is Hermitian Brownian motion $B(u)$ started at $u=1$ from a GUE matrix and observed at
$u=e^{2sn^{-1/3}}$, divided by $\sqrt u$, which is exactly the process $H$ of Section~\ref{sec:normal}. The tridiagonal
model uses the exact Gaussian transitions for the diagonal entries and the exact noncentral chi-square transitions for
$|Z_i|^2$. Neither model discretizes a stochastic differential equation. Ensemble replicates are independent; the four
observations within a replicate are coupled. Each pair of model and size draws its own stream from NumPy's
\texttt{SeedSequence} with base seed $4620261113$. Table~\ref{tab:sims} gives the results for the largest eigenvalue at
$s=1$, which is Airy time $1$.

\begin{table}[ht]
\centering
\caption{Largest edge-scaled eigenvalue at lags $0$ and $1$: correlation with $95\%$ percentile bootstrap interval
($500$ resamples) and mean squared increment $\EE(X_1(1)-X_1(0))^2$. The tridiagonal model has $2{,}000$ replicates at
every size.}\label{tab:sims}
\small
\begin{tabular}{rrcccc}
\toprule
& \multicolumn{3}{c}{correlation} & \multicolumn{2}{c}{mean squared increment}\\
\cmidrule(lr){2-4}\cmidrule(lr){5-6}
$n$ & dense replicates & dense & tridiagonal & dense & tridiagonal\\
\midrule
64 & 1500 & 0.430 [0.388, 0.476] & 0.631 [0.601, 0.658] & 0.972 & 0.589\\
128 & 1000 & 0.371 [0.315, 0.418] & 0.708 [0.687, 0.727] & 1.010 & 0.462\\
256 & 600 & 0.359 [0.277, 0.442] & 0.763 [0.746, 0.782] & 1.067 & 0.373\\
512 & 240 & 0.346 [0.229, 0.462] & 0.804 [0.788, 0.818] & 1.063 & 0.327\\
1024 & 64 & 0.524 [0.293, 0.685] & 0.842 [0.829, 0.854] & 0.863 & 0.251\\
2048 & 12 & 0.597 [0.120, 0.858] & 0.871 [0.861, 0.882] & 1.313 & 0.209\\
\bottomrule
\end{tabular}
\end{table}

The tridiagonal correlations increase and the tridiagonal increments decrease with $n$, as Theorem~\ref{thm:main}
predicts; the increments are still of order $0.2$ at $n=2048$, consistent with the slow rate in \eqref{eq:main}. The
dense rows at $n=1024$ and $n=2048$ have $64$ and $12$ replicates and wide intervals, and they do not establish a trend
for the dense model. The mean of $\Gamma(\lambda_1)$ under the Dumitriu--Edelman law falls from $0.361$ at $n=64$ to
$0.115$ at $n=2048$, against the value one for the dense process. A separate run with independent code and
$800$ tridiagonal replicates per size gave mean squared increments decreasing from $0.586$ at $n=64$ to $0.186$ at
$n=4096$, and a short-lag check at $n=20$, $\Delta s=10^{-4}$, gave dense increment variance per unit time
$1.936\pm0.035$ against the prediction $2$ from \eqref{eq:DBM}. The computations use NumPy~\cite{NumPy} and
SciPy~\cite{SciPy} and are not part of any proof.

\section*{Acknowledgements and funding}
This work was done by multiple AI agents, Claude (Anthropic) and Codex (OpenAI), connected to Math
Brain~\cite{MathBrainRepo} and working together inside the Brain framework~\cite{BrainFramework}, a shared, persistent
memory in which the agents record claims, proofs, failed attempts and methods, retrieve one another's results,
coordinate who works on what, and audit each other's work. For more on the details behind the Brain and Math Brain, see
the Math Brain paper~\cite{MathBrainPaper}, the Math Brain repository~\cite{MathBrainRepo} and the Brain framework
repository~\cite{BrainFramework}.

The research presented in this paper was supported by the European Research Council (ERC) under the European Union's
Horizon Europe research and innovation programme (grant agreement No.~101041711), by the Simons Foundation as part of
the Collaboration on the Mathematical and Scientific Foundations of Deep Learning, by Heights Labs, by Convex Nexus
Capital, by the Israel Science Foundation (grant number 2258/19), and by the Israel Science Foundation (ISF Grant
4101/25). Computational resources were provided in part by Google Cloud through its education credits programme (credit
530550103).

\section*{Declaration of competing interest}
The author holds equity in Heights Labs and in Convex Nexus Capital and has a management role at Convex Nexus Capital.
These entities are included in the funding declaration above.

\section*{Declaration of generative AI and AI-assisted technologies}
A Codex agent (OpenAI) proved the results and wrote the code of Section~\ref{sec:sims}; a second Codex agent checked the
proofs, the normalizations and the literature as a referee and ran the independent simulation; Claude agents (Anthropic)
checked the steps and prepared the manuscript. The author directed the work and is responsible for the content of the
paper and for decisions concerning its dissemination.

\section*{Code availability}
The code and data of Section~\ref{sec:sims} are in the repository
\begin{center}\url{https://github.com/yspennstate/tridiagonal-edge-freezing}.\end{center}

\appendix
\section{Notation}\label{app:notation}

\begin{center}
\begin{tabular}{p{3.0cm}p{8.0cm}p{3.1cm}}
\toprule
Symbol & Meaning & First appears\\
\midrule
$G$ & GUE matrix, density proportional to $\exp(-\Tr G^2/2)$ & Section~\ref{sec:process}\\
$a_i$, $b_i$ & diagonal and off-diagonal entries of a Jacobi matrix & \eqref{eq:DE}\\
$Z_i$ & vector Ornstein--Uhlenbeck process in $\RR^{2(n-i)}$, $b_i=|Z_i|/\sqrt2$ & \eqref{eq:process}\\
$T_n(t)$ & tridiagonal Ornstein--Uhlenbeck process & \eqref{eq:process}\\
$\lambda_{j,n}$, $v^{(j)}$ & $j$-th largest eigenvalue, its unit eigenvector & Section~\ref{sec:process}\\
$X_{j,n}(s)$ & edge-scaled eigenvalue at physical time $sn^{-1/3}$ & \eqref{eq:X}\\
$k$, $S$ & number of edge eigenvalues, time horizon & Theorem~\ref{thm:main}\\
$r_n$ & $\sqrt{12\log n}$ & Theorem~\ref{thm:main}\\
$N$ & dimension of the Gaussian process $Y$ & Section~\ref{sec:maximal}\\
$W$, $\widetilde W_i$, $B_j$ & standard Brownian motions & Sections~\ref{sec:process}--\ref{sec:maximal}\\
$\mathcal A_i$ & curves of the Airy line ensemble & Corollary~\ref{cor:frozen}\\
$H(t)$, $B(u)$ & Hermitian Ornstein--Uhlenbeck process, Hermitian Brownian motion & Section~\ref{sec:normal}\\
$Y$, $L$, $\gamma$, $\Gamma$ & Gaussian Ornstein--Uhlenbeck process, generator, invariant law, carr\'e du champ &
Section~\ref{sec:maximal}\\
$M$, $\widehat M$ & forward and backward martingales & Lemma~\ref{lem:maximal}\\
$E_n$ & event that all entries lie in their windows & \eqref{eq:En}\\
$m$, $A$, $B$ & block size $\lfloor n^{1/4}\rfloor$, leading block, Toeplitz comparison matrix & Lemma~\ref{lem:deloc}\\
$w_i$, $D$ & $|v^{(j)}_i|$, discrete Dirichlet energy of $w$ & Lemma~\ref{lem:deloc}\\
$c$, $\Gamma_n$, $K$ & entries $(a,b)$, carr\'e du champ, noise bound & Corollary~\ref{cor:bounded}\\
$C_n$ & covariance density of the entry noise & Proposition~\ref{prop:growth}\\
$\mathcal C$, $J$, $Q$, $K$ & cyclic set, Lanczos map, Lanczos basis, stabilizer of $e_1$ & Remark~\ref{rem:lanczos}\\
\bottomrule
\end{tabular}
\end{center}

\begin{thebibliography}{99}
\bibitem{AMP26} G.~Akemann, S.~N.~Majumdar and P.~P\"a{\ss}ler, A tridiagonal matrix-valued process with stochastic
resetting for arbitrary Dyson index $\beta>0$, preprint, 2026. \arxiv{2605.26945}
\bibitem{AG13} R.~Allez and A.~Guionnet, A diffusive matrix model for invariant $\beta$-ensembles, \emph{Electron.\ J.\
Probab.} \textbf{18} (2013), no.~62, 1--30. \doi{10.1214/EJP.v18-2073}
\bibitem{AGZ10} G.~W.~Anderson, A.~Guionnet and O.~Zeitouni, \emph{An Introduction to Random Matrices}, Cambridge
University Press, Cambridge, 2010. \doi{10.1017/CBO9780511801334}
\bibitem{CH14} I.~Corwin and A.~Hammond, Brownian Gibbs property for Airy line ensembles, \emph{Invent.\ Math.}
\textbf{195} (2014), 441--508. \doi{10.1007/s00222-013-0462-3}
\bibitem{DE02} I.~Dumitriu and A.~Edelman, Matrix models for beta ensembles, \emph{J.\ Math.\ Phys.} \textbf{43} (2002),
5830--5847. \doi{10.1063/1.1507823}
\bibitem{EJN24} A.~Edelman, S.~Jeong and R.~Nissim, On the limit of the tridiagonal model for $\beta$-Dyson Brownian
motion, preprint, 2024, revised 2026. \arxiv{2411.01633}
\bibitem{GXZ24} V.~Gorin, J.~Xu and L.~Zhang, Airy$_\beta$ line ensemble and its Laplace transform, preprint, 2024.
\arxiv{2411.10829}
\bibitem{NumPy} C.~R.~Harris et al., Array programming with NumPy, \emph{Nature} \textbf{585} (2020), 357--362.
\doi{10.1038/s41586-020-2649-2}
\bibitem{HP17} D.~Holcomb and E.~Paquette, Tridiagonal models for Dyson Brownian motion, preprint, 2017.
\arxiv{1707.02700}
\bibitem{HZ26} J.~Huang and L.~Zhang, A convergence framework for Airy$_\beta$ line ensemble via pole evolution,
\emph{Adv.\ Math.} \textbf{498} (2026), 111028. \doi{10.1016/j.aim.2026.111028}
\bibitem{Lan20} B.~Landon, Edge scaling limit of Dyson Brownian motion at equilibrium for general $\beta\ge1$, preprint,
2020. \arxiv{2009.11176}
\bibitem{LM00} B.~Laurent and P.~Massart, Adaptive estimation of a quadratic functional by model selection, \emph{Ann.\
Statist.} \textbf{28} (2000), 1302--1338. \doi{10.1214/aos/1015957395}
\bibitem{Li10} L.-C.~Li, The beta-Hermite and beta-Laguerre processes, preprint, 2010. \arxiv{1007.3905}
\bibitem{LZ88} T.~Lyons and W.~Zheng, A crossing estimate for the canonical process on a Dirichlet space and a
tightness result, in \emph{Colloque Paul L\'evy sur les processus stochastiques}, Ast\'erisque \textbf{157--158}
(1988), 249--271, Soci\'et\'e Math\'ematique de France, Paris.
\bibitem{RRV11} J.~A.~Ram\'irez, B.~Rider and B.~Vir\'ag, Beta ensembles, stochastic Airy spectrum, and a diffusion,
\emph{J.\ Amer.\ Math.\ Soc.} \textbf{24} (2011), 919--944. \doi{10.1090/S0894-0347-2011-00703-0}
\bibitem{MathBrainPaper} Y.~Shmalo, \emph{Math Brain: shared, auditable memory for multi-agent mathematical research},
October 2026, \url{https://github.com/yspennstate/math-brain-public/tree/main/papers/math-brain-2026-10-24}.
\bibitem{MathBrainRepo} Y.~Shmalo, Math Brain: public snapshot, code and tests, 2026,
\url{https://github.com/yspennstate/math-brain-public}.
\bibitem{BrainFramework} Y.~Shmalo, Brain Framework, public software repository, 2026,
\url{https://github.com/yspennstate/brain-framework-public}.
\bibitem{TW94} C.~A.~Tracy and H.~Widom, Level-spacing distributions and the Airy kernel, \emph{Comm.\ Math.\ Phys.}
\textbf{159} (1994), 151--174. \doi{10.1007/BF02100489}
\bibitem{VV17} B.~Valk\'o and B.~Vir\'ag, The Sine$_\beta$ operator, \emph{Invent.\ Math.} \textbf{209} (2017),
275--327. \doi{10.1007/s00222-016-0709-x}
\bibitem{Vir14} B.~Vir\'ag, Operator limits of random matrices, in \emph{Proceedings of the International Congress of
Mathematicians, Seoul 2014}, Vol.~4, 247--272. \arxiv{1804.06953}
\bibitem{SciPy} P.~Virtanen et al., SciPy 1.0: fundamental algorithms for scientific computing in Python, \emph{Nature
Methods} \textbf{17} (2020), 261--272. \doi{10.1038/s41592-019-0686-2}
\bibitem{Yab23} S.~Yabuoku, Eigenvalue processes of symmetric tridiagonal matrix-valued processes associated with
Gaussian beta ensemble, \emph{Stochastic Process.\ Appl.} \textbf{164} (2023), 139--159. \doi{10.1016/j.spa.2023.07.004}
\end{thebibliography}
\end{document}
